\section*{Introduction}
\sectionmark{Introduction}


\begin{quote}
    \textit{Dans l'état actuel des connaissances en sémiologie, l'image apparaît comme un univers à la complexité redoutable : ses lois sont mal connues. Cela par opposition à la parole linguistique, qui semble plus accessible, mieux apprivoisée... Or l'une et l'autre apparaissent souvent ensemble [...] et parfois étroitement mêlées}\vspace{0.2cm}
    
    Laurence Badin, \og  Le texte et l'image \fg, dans \textit{Communication \& Langages}, vol 26, pp. 98-112 , 1975
\end{quote}

\lettrine{L}{'application de méthodes} informatiques sur les sujets de sciences humaines et sociales a toujours une forme de bricolage. Comme souvent, le chercheur se retrouve à utiliser des outils qui n'ont pas été pensés pour sa discipline et se doit de les réadapter pour sa pratique scientifique ; le domaine de la vision artificielle ne fait pas exception. Dans cette partie, nous nous intéresserons aux modes de développement et au fonctionnement de l'écosystème international des technologies que nous essayons de réadapter à des questions historiques. 

D'abord, nous contextualiserons ce qu'est la vision artificielle et les différentes phases techniques et conceptuelles par lesquelles la discipline s'est formée pour en arriver aux technologies actuelles applicables à des documents historiques. Ensuite, nous étudierons les différents tests, déploiements et expériences de l'utilisation de la vision artificielle dans le domaine de l'étude du document et des institutions patrimoniales. Enfin, nous nous concentrerons sur une réflexion historienne sur la place de ces outils en fonction des différents besoins de recherche. 